\chapter{Painted Bipartitions (PBP)}
\label{chap:pbp}

Painted Bipartitions are the basic combinatorial data parametrising local
systems on nilpotent orbits (they encode both the $\wp$-decorations and the
local system type). We adopt the definitions of [BMSZb] Section~2.

\section{DRC symbols and layer order}

\begin{definition}[DRC symbols]
  \label{def:drc_symbols}
  \lean{DRCSymbol}
  \leanok
  The five \emph{DRC symbols} are
  \[ \{\bullet,\ s,\ r,\ c,\ d\}. \]
  Symbols other than $\bullet$ are called \emph{non-dot}.
\end{definition}

\begin{definition}[Layer order]
  \label{def:layerOrd}
  \lean{DRCSymbol.layerOrd}
  \leanok
  \uses{def:drc_symbols}
  The \emph{layer order} is a function $\mathrm{layerOrd} : \mathrm{DRCSymbol} \to \mathbb{N}$ given by
  \[ \mathrm{layerOrd}(\bullet) = 0,\ \mathrm{layerOrd}(s) = 1,\ \mathrm{layerOrd}(r) = 2,\ \mathrm{layerOrd}(c) = 3,\ \mathrm{layerOrd}(d) = 4. \]
  This is the linear order $\bullet < s < r < c < d$.
\end{definition}

\begin{definition}[Tail contribution]
  \label{def:tailContrib}
  \lean{DRCSymbol.tailContrib}
  \leanok
  \uses{def:drc_symbols}
  Paper Eq.~(11.5). Each DRC symbol contributes a pair of non-negative integers
  to the tail signature:
  \[
    \mathrm{tailContrib}(\bullet) = (1,1),\ \mathrm{tailContrib}(s) = (0, 2),\ \mathrm{tailContrib}(r) = (2, 0),
  \]
  \[
    \mathrm{tailContrib}(c) = (1, 1),\ \mathrm{tailContrib}(d) = (1, 1).
  \]
  In particular the two components always sum to~$2$.
\end{definition}

\section{Painted Young diagrams}

\begin{definition}[Young diagram]
  \label{def:yd}
  % YoungDiagram from Mathlib
  \leanok
  A Young diagram $\mu$ is a finite subset of $\mathbb{N}^2$ which is a
  \emph{down-set} under $(i, j) \leq (i', j') \iff i \leq i' \land j \leq j'$.
  Columns $j$ have lengths $\mu.\mathrm{colLen}(j) = |\{i : (i, j) \in \mu\}|$.
\end{definition}

\begin{definition}[Painted Young diagram]
  \label{def:pyd}
  \lean{PaintedYoungDiagram.layerMonotone}
  \leanok
  \uses{def:drc_symbols, def:layerOrd, def:yd}
  A \emph{painted Young diagram} is a triple $D = (\mu, \mathcal{P}, \mathrm{layerMono})$ where:
  \begin{enumerate}
    \item $\mu$ is a Young diagram (the \emph{shape}).
    \item $\mathcal{P} : \mathbb{N}^2 \to \mathrm{DRCSymbol}$ is a paint function with
      $\mathcal{P}(i, j) = \bullet$ for $(i, j) \notin \mu$ (paint outside the shape is $\bullet$).
    \item \textbf{layer monotonicity}: for all $(i_1, j_1), (i_2, j_2)$ with $i_1 \leq i_2$ and $j_1 \leq j_2$,
      if $(i_2, j_2) \in \mu$ then $\mathrm{layerOrd}(\mathcal{P}(i_1, j_1)) \leq \mathrm{layerOrd}(\mathcal{P}(i_2, j_2))$.
  \end{enumerate}
\end{definition}

\begin{definition}[$\mathrm{dotScolLen}$ — dot-and-s column length]
  \label{def:dotScolLen}
  \lean{PBP.dotScolLen}
  \leanok
  \uses{def:layerOrd, def:pyd}
  For a painted Young diagram $D$, $\mathrm{dotScolLen}(D, j)$ is the number of cells in column $j$ whose paint is $\bullet$ or $s$ (equivalently, the largest $k$ such that every cell $(i, j)$ with $i < k$ has $\mathrm{layerOrd}(\mathcal{P}(i, j)) \leq 1$).
\end{definition}

\section{Root type and symbol allowances}

\begin{definition}[Root type]
  \label{def:RootType}
  \lean{RootType}
  \leanok
  A \emph{root type} $\gamma \in \{B^+,\ B^-,\ C,\ D,\ M\}$ selects a set of allowed paint symbols for $\mathcal{P}$ and $\mathcal{Q}$ (see below).
\end{definition}

\begin{definition}[Symbol allowances per root type]
  \label{def:sym_allowed}
  \lean{DRCSymbol.allowed}
  \leanok
  \uses{def:RootType, def:drc_symbols}
  The following table gives the allowed DRC symbols for $\mathcal{P}$ and $\mathcal{Q}$ (paper Def.~2.25):
  \[
    \begin{array}{l|ll}
      \gamma & \mathcal{P}\text{ allowed} & \mathcal{Q}\text{ allowed} \\\hline
      B^+ / B^- & \{\bullet, c\} & \{\bullet, s, r, d\} \\
      C        & \{\bullet, r, c, d\} & \{\bullet, s\} \\
      D        & \{\bullet, s, r, c, d\} & \{\bullet\} \\
      M = \tilde{C}        & \{\bullet, s, c\} & \{\bullet, r, d\} \\
    \end{array}
  \]
  The $B^+/B^-$ distinction corresponds to a tag on the extended DRC diagram.
\end{definition}

\section{Painted Bipartition}

\begin{definition}[Painted Bipartition]
  \label{def:PBP}
  \lean{PBP}
  \leanok
  \uses{def:RootType, def:drc_symbols, def:pyd, def:sym_allowed}
  A \emph{painted bipartition} (PBP) is $\tau = (\mathcal{P}, \mathcal{Q}, \gamma)$ where
  $\mathcal{P}$ and $\mathcal{Q}$ are painted Young diagrams with root type $\gamma$
  satisfying:
  \begin{enumerate}
    \item \textbf{sym\_P, sym\_Q}: $\mathcal{P}.\mathrm{paint}(i, j)$ is in the allowed set for $\mathcal{P}$ determined by $\gamma$ (Def.~\ref{def:sym_allowed}), and likewise for $\mathcal{Q}$.
    \item \textbf{dot\_match}: $\mathcal{P}(i,j) = \bullet \Leftrightarrow \mathcal{Q}(i,j) = \bullet$ (the outside-shape convention makes this an equivalence between cell memberships).
    \item \textbf{row\_s, row\_r}: each row contains at most one $s$ across $\mathcal{P} \cup \mathcal{Q}$; similarly at most one $r$.
    \item \textbf{col\_c\_P / col\_c\_Q}: each column of $\mathcal{P}$ has at most one $c$; same for $\mathcal{Q}$.
    \item \textbf{col\_d\_P / col\_d\_Q}: each column has at most one $d$.
  \end{enumerate}
\end{definition}

\begin{remark}
  The notation $\iota_\wp := \mathcal{P}$ and $\jmath_\wp := \mathcal{Q}$ appears in the paper and is used interchangeably. The index $\wp$ is a \emph{primitive pair set} which we usually omit, assuming the canonical $\wp$.
\end{remark}

\begin{definition}[PBPSet — PBPs of fixed shape]
  \label{def:PBPSet}
  \lean{PBPSet}
  \leanok
  \uses{def:PBP}
  For $\gamma$ a root type and $\mu_P, \mu_Q$ Young diagrams,
  \[ \PBP_\gamma(\mu_P, \mu_Q) := \{\tau \in \mathrm{PBP} \mid \tau.\gamma = \gamma,\ \tau.\mathcal{P}.\mathrm{shape} = \mu_P,\ \tau.\mathcal{Q}.\mathrm{shape} = \mu_Q\}. \]
  Given a dual partition $dp$, we set $\PBP_\gamma(\check{\mathcal{O}}) := \PBP_\gamma(\mu_P(dp), \mu_Q(dp))$ where $\mu_P, \mu_Q$ are derived from $dp$.
\end{definition}

\section{Signature and epsilon}

\begin{definition}[PBP signature $\Sign(\tau)$]
  \label{def:signature_pbp}
  \lean{PBP.signature}
  \leanok
  \uses{def:PBP}
  The signature $\Sign(\tau) = (p_\tau, q_\tau) \in \mathbb{N}^2$ counts paint
  symbols weighted by paint multiplicities (paper Eq.~(3.3)). Specifically,
  denoting by $\#_\sigma$ the count of symbol $\sigma$ in $\mathcal{P}$ and $\mathcal{Q}$,
  and by $n_s, n_r$ the total $s, r$ counts across both,
  \[
    p_\tau = n_r, \qquad q_\tau = n_s,
  \]
  with appropriate bookkeeping for $\bullet, c, d$ counts and per-$\gamma$ variations.
  For D type: $p_\tau = \#_{\bullet,\iota} + \#_{\bullet,\jmath} + 2\#_{r,\iota} + \#_{c,\iota} + \#_{d,\iota}$ (and similarly for $q_\tau$ with $s$ in place of $r$).
\end{definition}

\begin{definition}[Epsilon $\eps_\tau$]
  \label{def:epsilon_pbp}
  \lean{PBP.epsilon}
  \leanok
  \uses{def:PBP}
  (Paper Eq.~(3.6).) $\eps_\tau \in \{0, 1\}$ is $1$ iff column~0 of $\mathcal{P}$ (or $\mathcal{Q}$, per $\gamma$) contains a $d$.
\end{definition}

\begin{definition}[Total cell count $|\tau|$]
  \label{def:cardCells}
  \lean{PBP.cardCells}
  \leanok
  \uses{def:PBP}
  $|\tau| := |\mathcal{P}.\mathrm{shape}.\mathrm{cells}| + |\mathcal{Q}.\mathrm{shape}.\mathrm{cells}|$.
\end{definition}

\section{Dual partitions and shape derivation}

\begin{definition}[Dual partition legality]
  \label{def:DualPart}
  \lean{DualPart, }
  \leanok
  A \emph{dual partition} $dp$ is a \texttt{List ℕ} that is sorted $\geq$ (strictly decreasing with possible ties). Validity conditions:
  \begin{itemize}
    \item B-type / M-type: every entry is even and positive.
    \item D-type / C-type: every entry is odd.
  \end{itemize}
\end{definition}

\begin{definition}[Column lengths from dp]
  \label{def:dpartColLens}
  \lean{dpartColLensP_B, dpartColLensQ_B, dpartColLensP_D, dpartColLensQ_D, dpartColLensP_C, dpartColLensQ_C, dpartColLensP_M, dpartColLensQ_M}
  \leanok
  \uses{def:DualPart}
  For each root type, functions compute the column-length lists of $\mu_P$ and $\mu_Q$ from $dp$. For example (D type): $\mu_P(dp)$ has column lengths $(dp_1 + 1)/2, (dp_2 + 1)/2, \ldots$, etc.
\end{definition}
